\ifdefined\BookMaster
\def\BookEnd{}
\else
\documentclass[a4paper,12pt,fleqn,openany,twoside,fontset=windows]{ctexbook}

\usepackage[fleqn]{amsmath}
\usepackage{amssymb,amsthm,mathrsfs}
\usepackage{geometry}
\usepackage[dvipsnames]{xcolor}
\usepackage{graphicx}
\usepackage{tikz}
\usetikzlibrary{calc}
\usepackage[most]{tcolorbox}
\usepackage{fancyhdr}
\usepackage{titlesec}
\usepackage{tocloft}
\usepackage{enumitem}
\usepackage{microtype}
\usepackage{pifont}
\usepackage{hyperref}
\usepackage{romannum}
\usepackage{mathtools}

\definecolor{bookblue}{HTML}{264653}
\definecolor{bookteal}{HTML}{4F7C82}
\definecolor{booklight}{HTML}{EDF4F3}
\definecolor{bookgray}{HTML}{5E686B}

\geometry{
    inner=2.7cm,
    outer=2.5cm,
    top=2.7cm,
    bottom=2.7cm,
    headheight=18pt,
    headsep=18pt,
    footskip=25pt
}
\setlength{\mathindent}{0pt}
\setlength{\parindent}{5pt}
\setlength{\parskip}{3pt plus 1pt minus 1pt}
\linespread{1.25}
\allowdisplaybreaks[3]
\AtBeginDocument{%
    \setlength{\parindent}{5pt}
    \setlength{\abovedisplayskip}{8pt plus 2pt minus 2pt}
    \setlength{\belowdisplayskip}{8pt plus 2pt minus 2pt}
    \setlength{\abovedisplayshortskip}{5pt plus 1pt minus 1pt}
    \setlength{\belowdisplayshortskip}{5pt plus 1pt minus 1pt}
}

\hypersetup{
    unicode=true,
    colorlinks=true,
    linkcolor=bookblue,
    citecolor=bookblue,
    urlcolor=bookteal,
    bookmarksopen=true,
    pdfauthor={Chen},
    pdftitle={数学分析习题整理}
}

\renewcommand{\chaptermark}[1]{%
    \edef\@tmp{第\thechapter 章\quad #1}%
    \expandafter\markboth\expandafter{\@tmp}{}%
}
\fancypagestyle{main}{%
    \fancyhf{}
    \fancyhead[LE]{\small\color{bookgray}\nouppercase{\leftmark}}
    \fancyhead[RO]{\small\color{bookgray}数学分析习题整理}
    \fancyfoot[C]{\color{bookblue}\thepage}
    \fancyfoot[R]{\small\bfseries Chen\;\; github.com/mathlover-1/math-analysis}
    \renewcommand{\headrulewidth}{0.4pt}
    \renewcommand{\footrulewidth}{0pt}
}
\fancypagestyle{plain}{%
    \fancyhf{}
    \fancyfoot[C]{\color{bookblue}\thepage}
    \fancyfoot[R]{\small\bfseries Chen\;\; github.com/mathlover-1/math-analysis}
    \renewcommand{\headrulewidth}{0pt}
    \renewcommand{\footrulewidth}{0pt}
}
\pagestyle{main}

\titleformat{\chapter}[display]
  {\normalfont\sffamily\bfseries\color{bookblue}}
  {\raggedleft\fontsize{46}{50}\selectfont\thechapter}
  {-18pt}
  {\titlerule[1.2pt]\vspace{10pt}\filright\Huge}
\titlespacing*{\chapter}{0pt}{-15pt}{36pt}

\renewcommand{\contentsname}{目\quad 录}
\renewcommand{\cfttoctitlefont}{\hfill\sffamily\bfseries\Huge\color{bookblue}}
\renewcommand{\cftaftertoctitle}{\hfill}
\setlength{\cftbeforetoctitleskip}{5pt}
\setlength{\cftaftertoctitleskip}{28pt}
\setlength{\cftbeforechapskip}{12pt}
\renewcommand{\cftchapfont}{\sffamily\bfseries\large\color{bookblue}}
\renewcommand{\cftchappagefont}{\sffamily\bfseries\large\color{bookblue}}
\renewcommand{\cftchapleader}{\color{bookteal}\cftdotfill{2.5}}

\newtcolorbox{ti}{
    enhanced,
    breakable,
    colback=booklight!55!white,
    colframe=bookteal!35!white,
    boxrule=0.45pt,
    boxsep=0pt,
    left=10pt,
    right=10pt,
    top=9pt,
    bottom=9pt,
    before skip=14pt,
    after skip=14pt,
    arc=2pt,
    outer arc=2pt,
    before upper={\parindent=0pt}
}

\newenvironment{booknotebody}{%
    \begin{center}
    \begin{minipage}{0.84\textwidth}
    \songti\zihao{-4}\linespread{1.45}\selectfont
    \setlength{\parindent}{2\ccwd}
    \setlength{\parskip}{0.75em}
}{%
    \end{minipage}
    \end{center}
}

\newenvironment{booknotelongbody}{%
    \begingroup
    \songti\zihao{-4}\linespread{1.45}\selectfont
    \setlength{\parindent}{2\ccwd}
    \setlength{\parskip}{0.75em}
    \begin{list}{}{%
        \setlength{\leftmargin}{0.08\textwidth}
        \setlength{\rightmargin}{0.08\textwidth}
        \setlength{\topsep}{0pt}
        \setlength{\partopsep}{0pt}
        \setlength{\parsep}{0pt}
        \setlength{\itemsep}{0pt}
    }
    \item\relax
}{%
    \end{list}
    \endgroup
}

\fancypagestyle{plainnowm}{%
    \fancyhf{}
    \fancyfoot[C]{\color{bookblue}\thepage}
    \renewcommand{\headrulewidth}{0pt}
    \renewcommand{\footrulewidth}{0pt}
}

\newcommand{\booknotetitle}[2][plain]{%
    \clearpage
    \phantomsection
    \addcontentsline{toc}{chapter}{#2}
    \markboth{#2}{}
    \thispagestyle{#1}
    \vspace*{1.2cm}
    \begin{center}
        {\sffamily\heiti\bfseries\fontsize{26}{34}\selectfont\color{bookblue}#2\par}
        \vspace{0.8em}
        {\color{bookteal}\rule{4.5cm}{0.8pt}\par}
    \end{center}
    \vspace{0.6cm}
}

\renewcommand{\proofname}{证明}
\renewcommand{\qedsymbol}{\color{bookteal}\ensuremath{\blacksquare}}

\title{数学分析习题整理}
\author{Chen}
\date{}

\makeatletter
% ==================== 旧版封面/封底设计备份 ====================
% 2026-08 重设计,旧代码用 \iffalse...\fi 封存,新版见下方。
\iffalse
\renewcommand{\maketitle}{%
    \begin{titlepage}
        \thispagestyle{empty}
        \setcounter{page}{0}
        \noindent\hspace*{-\parindent}\begin{tikzpicture}[remember picture,overlay]
            \fill[bookblue] (current page.north west) rectangle ([yshift=-4.2cm]current page.north east);
            \fill[booklight] ([yshift=-4.2cm]current page.north west) rectangle (current page.south east);
            \draw[bookteal!55,line width=0.7pt]
                ([xshift=2.4cm,yshift=4.3cm]current page.south west)
                .. controls +(2.5cm,1.8cm) and +(-2.8cm,-1.3cm) ..
                ([xshift=-2.4cm,yshift=7.2cm]current page.south east);
            \draw[bookteal!35,line width=0.7pt]
                ([xshift=2.5cm,yshift=6.2cm]current page.south west)
                .. controls +(3.0cm,-1.2cm) and +(-3.0cm,1.6cm) ..
                ([xshift=-2.3cm,yshift=3.8cm]current page.south east);
        \end{tikzpicture}
        \vspace*{5.6cm}
        \noindent\color{bookblue}\rule{\textwidth}{1.4pt}\par
        \vspace{8pt}
        {\centering\sffamily\bfseries\fontsize{32}{42}\selectfont\color{bookblue}\@title\par}
        \vspace{8pt}
        \noindent\color{bookteal}\rule{\textwidth}{0.6pt}\par
        \vspace{2.6cm}
        {\centering\Large\color{bookgray}\@author\par}
        \vfill
        {\centering\color{bookteal}\large
            $\displaystyle \lim_{n\to\infty} a_n$\qquad
            $\displaystyle \int_a^b f(x)\,\mathrm{d}x$\qquad
            $\displaystyle \nabla f(x,y)$\par
        }
        \vspace*{1.5cm}
    \end{titlepage}
}

\newcommand{\makebackcover}{%
    \begingroup
    \let\cleardoublepage\clearpage
    \begin{titlepage}
        \thispagestyle{empty}
        \noindent\hspace*{-\parindent}\begin{tikzpicture}[remember picture,overlay]
            \fill[bookblue] (current page.north west) rectangle (current page.south east);
            \fill[booklight,rounded corners=7pt]
                ([xshift=2.2cm,yshift=-3.0cm]current page.north west)
                rectangle
                ([xshift=-2.2cm,yshift=3.0cm]current page.south east);
            \fill[bookteal!80]
                (current page.north west) --
                ([xshift=5.3cm]current page.north west) --
                ([xshift=2.1cm,yshift=-4.4cm]current page.north west) --
                ([yshift=-6.2cm]current page.north west) -- cycle;
            \fill[bookteal!55]
                (current page.south east) --
                ([xshift=-5.8cm]current page.south east) --
                ([xshift=-2.4cm,yshift=4.2cm]current page.south east) --
                ([yshift=6.0cm]current page.south east) -- cycle;
            \foreach \radius in {0.75,1.15,1.55,1.95} {
                \draw[bookteal!55,line width=0.55pt]
                    ([xshift=-2.9cm,yshift=-4.4cm]current page.north east)
                    circle[radius=\radius cm];
            }
            \foreach \radius in {0.65,1.05,1.45} {
                \draw[booklight!65,line width=0.55pt]
                    ([xshift=2.4cm,yshift=3.6cm]current page.south west)
                    circle[radius=\radius cm];
            }
            \fill[bookteal!16,rounded corners=3pt]
                ([xshift=3.5cm,yshift=-7.1cm]current page.north west)
                rectangle
                ([xshift=-5.0cm,yshift=7.2cm]current page.south east);
            \fill[bookteal!30,rounded corners=3pt]
                ([xshift=4.4cm,yshift=-8.0cm]current page.north west)
                rectangle
                ([xshift=-4.1cm,yshift=8.1cm]current page.south east);
            \begin{scope}[shift={([xshift=-0.75cm]current page.center)}]
                \clip[rounded corners=3pt] (-5.25,-4.75) rectangle (5.25,4.75);
                \foreach \v in {-3.6,-3.0,...,3.6} {
                    \draw[bookteal!48,line width=0.45pt,opacity=0.72]
                        plot[domain=-3.6:3.6,samples=45,smooth,variable=\u]
                        ({\u+0.55*\v},{0.28*\u-0.35*\v+0.11*(\u*\u-\v*\v)});
                }
                \foreach \u in {-3.6,-3.0,...,3.6} {
                    \draw[bookblue!42,line width=0.42pt,opacity=0.62]
                        plot[domain=-3.6:3.6,samples=45,smooth,variable=\v]
                        ({\u+0.55*\v},{0.28*\u-0.35*\v+0.11*(\u*\u-\v*\v)});
                }
                \draw[bookteal!70,line width=0.75pt,opacity=0.82]
                    plot[domain=-3.6:3.6,samples=55,smooth,variable=\u]
                    ({\u},{0.28*\u+0.11*\u*\u});
                \draw[bookblue!60,line width=0.7pt,opacity=0.75]
                    plot[domain=-3.6:3.6,samples=55,smooth,variable=\v]
                    ({0.55*\v},{-0.35*\v-0.11*\v*\v});
            \end{scope}
            \draw[bookteal!45,line width=0.65pt]
                ([xshift=3.0cm,yshift=-8.2cm]current page.north west)
                .. controls +(2.2cm,2.0cm) and +(-2.0cm,-2.1cm) ..
                ([xshift=-3.0cm,yshift=9.0cm]current page.south east);
            \draw[bookteal!25,line width=0.55pt]
                ([xshift=3.1cm,yshift=8.0cm]current page.south west)
                .. controls +(2.7cm,-1.7cm) and +(-2.6cm,1.8cm) ..
                ([xshift=-3.1cm,yshift=-8.7cm]current page.north east);
            \node[anchor=west,text=bookteal!72,scale=0.78] at
                ([xshift=2.8cm,yshift=-4.0cm]current page.north west) {%
                $\displaystyle e^{\mathrm{i}\pi}+1=0$};
            \node[anchor=east,text=bookblue!62,scale=0.67] at
                ([xshift=-2.7cm,yshift=4.0cm]current page.south east) {%
                $\displaystyle
                f(x)=\sum_{n=0}^{\infty}\frac{f^{(n)}(a)}{n!}(x-a)^n$};
        \end{tikzpicture}
        \mbox{}
    \end{titlepage}
    \endgroup
}
\fi

% ==================== 新版封面(2026-08 重设计) ====================
% 设计要点:顶部深绿纵向渐变+波浪下缘+两层半透明过渡波(解决深浅绿生硬过渡);
% 深色区螺旋线/同心圆母题;标题+英文副题+作者居中;
% 中部 Fourier 部分和逼近方波曲线;底部浅色波纹(无文字)。
\renewcommand{\maketitle}{%
    \begin{titlepage}
        \thispagestyle{empty}
        \setcounter{page}{0}
        \noindent\hspace*{-\parindent}\begin{tikzpicture}[remember picture,overlay]
            % ---------- 底色 ----------
            \fill[booklight] (current page.north west) rectangle (current page.south east);

            % ---------- 顶部深色区:纵向渐变 + 波浪下缘 ----------
            \shade[top color=bookblue, bottom color=bookteal!55!bookblue]
                (current page.north west) -- (current page.north east) --
                ([yshift=-8.6cm]current page.north east)
                .. controls ([xshift=-5.2cm,yshift=-10.6cm]current page.north east)
                         and ([xshift=6.0cm,yshift=-8.0cm]current page.north west) ..
                ([yshift=-10.0cm]current page.north west) -- cycle;

            % ---------- 过渡层 1:半透明波浪 ----------
            \fill[bookteal!60!bookblue,opacity=0.42]
                ([yshift=-8.6cm]current page.north east)
                .. controls ([xshift=-5.2cm,yshift=-10.6cm]current page.north east)
                         and ([xshift=6.0cm,yshift=-8.0cm]current page.north west) ..
                ([yshift=-10.0cm]current page.north west) --
                ([yshift=-12.6cm]current page.north west)
                .. controls ([xshift=6.5cm,yshift=-11.0cm]current page.north west)
                         and ([xshift=-5.5cm,yshift=-13.2cm]current page.north east) ..
                ([yshift=-11.4cm]current page.north east) -- cycle;

            % ---------- 过渡层 2:更浅的波浪 ----------
            \fill[bookteal!35,opacity=0.35]
                ([yshift=-10.6cm]current page.north east)
                .. controls ([xshift=-6.0cm,yshift=-12.4cm]current page.north east)
                         and ([xshift=5.0cm,yshift=-9.8cm]current page.north west) ..
                ([yshift=-11.8cm]current page.north west) --
                ([yshift=-13.8cm]current page.north west)
                .. controls ([xshift=5.5cm,yshift=-12.6cm]current page.north west)
                         and ([xshift=-6.5cm,yshift=-14.6cm]current page.north east) ..
                ([yshift=-12.8cm]current page.north east) -- cycle;

            % ---------- 深色区装饰:右上螺旋线 ----------
            \begin{scope}[shift={([xshift=-3.4cm,yshift=-4.3cm]current page.north east)}]
                \draw[booklight!45!bookteal,opacity=0.55,line width=0.6pt]
                    plot[domain=0:1440,samples=220,smooth,variable=\t]
                    ({0.0028*\t*cos(\t)},{0.0028*\t*sin(\t)});
            \end{scope}

            % ---------- 深色区装饰:左下同心圆 ----------
            \foreach \r in {0.7,1.25,1.8,2.35,2.9}{
                \draw[booklight!35!bookteal,opacity=0.5,line width=0.5pt]
                    ([xshift=0.4cm,yshift=-7.9cm]current page.north west) circle[radius=\r cm];
            }

            % ---------- 主标题 ----------
            \node[anchor=center,text=booklight] at
                ([yshift=-4.5cm]current page.north)
                {\sffamily\bfseries\fontsize{40}{46}\selectfont \@title};

            % ---------- 分隔短线 ----------
            \draw[booklight!65,line width=1.2pt]
                ([xshift=-2.6cm,yshift=-5.75cm]current page.north) --
                ([xshift=2.6cm,yshift=-5.75cm]current page.north);

            % ---------- 英文副题 ----------
            \node[anchor=center,text=booklight!70!bookteal] at
                ([yshift=-6.55cm]current page.north)
                {\itshape\large Mathematical Analysis \textperiodcentered\
                 A Collection of Exercises};

            % ---------- 作者(紧随标题下方) ----------
            \node[anchor=center,text=booklight!88] at
                ([yshift=-8.2cm]current page.north)
                {\sffamily\large \@author\quad 整理};

            % ---------- 中部:Fourier 部分和逼近方波 ----------
            \begin{scope}[shift={([xshift=-5.2cm,yshift=-18.1cm]current page.north)}]
                % 坐标轴
                \draw[bookgray!55,line width=0.5pt,->] (-0.5,0) -- (10.9,0);
                \draw[bookgray!55,line width=0.5pt,->] (0,-1.75) -- (0,1.95);
                % 目标方波(虚线)
                \draw[bookgray!75,line width=0.6pt,dashed]
                    (0,0.92) -- (5,0.92) (5,-0.92) -- (10,-0.92);
                % N=1
                \draw[bookteal!45,line width=0.7pt]
                    plot[domain=0:360,samples=90,smooth,variable=\x]
                    ({\x/36},{1.18*sin(\x)});
                % N=3
                \draw[bookteal!80,line width=0.8pt]
                    plot[domain=0:360,samples=120,smooth,variable=\x]
                    ({\x/36},{1.18*(sin(\x)+sin(3*\x)/3)});
                % N=7
                \draw[bookblue,line width=0.95pt]
                    plot[domain=0:360,samples=150,smooth,variable=\x]
                    ({\x/36},{1.18*(sin(\x)+sin(3*\x)/3+sin(5*\x)/5+sin(7*\x)/7)});
                % 刻度标注
                \node[anchor=north east,text=bookgray!80,scale=0.62] at (-0.12,0) {$0$};
                \node[anchor=north,text=bookgray!80,scale=0.62] at (5,0) {$\pi$};
                \node[anchor=north,text=bookgray!80,scale=0.62] at (10,0) {$2\pi$};
            \end{scope}
            % 图注
            \node[anchor=center,text=bookgray,scale=0.8] at
                ([yshift=-20.5cm]current page.north)
                {$S_N(x)=\displaystyle\sum_{k=1}^{N}\frac{\sin\bigl((2k-1)x\bigr)}{2k-1}$
                 \ 逐步逼近方波};

            % ---------- 底部浅色波纹(无文字,与顶部波浪呼应) ----------
            \fill[bookteal!22,opacity=0.5]
                ([yshift=2.4cm]current page.south west)
                .. controls ([xshift=6.0cm,yshift=3.9cm]current page.south west)
                         and ([xshift=-6.5cm,yshift=1.0cm]current page.south east) ..
                ([yshift=2.9cm]current page.south east) --
                (current page.south east) -- (current page.south west) -- cycle;
            \fill[bookteal!40!bookblue,opacity=0.28]
                ([yshift=1.3cm]current page.south west)
                .. controls ([xshift=5.5cm,yshift=2.7cm]current page.south west)
                         and ([xshift=-5.5cm,yshift=0.3cm]current page.south east) ..
                ([yshift=1.9cm]current page.south east) --
                (current page.south east) -- (current page.south west) -- cycle;
        \end{tikzpicture}
        \mbox{}
    \end{titlepage}
}

% ==================== 新版封底(纯背景,2026-08 重设计) ====================
% 设计要点:对角渐变底+角部半透明叠层;中下部外摆线(玫瑰形)为视觉主体;
% 底部两层半透明白色波浪;无任何文字。
\newcommand{\makebackcover}{%
    \begingroup
    \let\cleardoublepage\clearpage
    \begin{titlepage}
        \thispagestyle{empty}
        \noindent\hspace*{-\parindent}\begin{tikzpicture}[remember picture,overlay]
            % ---------- 底色:对角渐变 ----------
            \shade[top color=bookblue, bottom color=bookteal!48!bookblue,
                   shading angle=-30]
                (current page.north west) rectangle (current page.south east);

            % ---------- 角部半透明叠层 ----------
            \fill[white,opacity=0.05]
                ([xshift=2.0cm,yshift=1.0cm]current page.north east) circle[radius=7.2cm];
            \fill[white,opacity=0.06]
                ([xshift=-0.5cm,yshift=-1.5cm]current page.north east) circle[radius=3.6cm];

            % ---------- 左下同心圆 ----------
            \foreach \r in {1.0,1.75,2.5,3.25,4.0,4.75}{
                \draw[booklight!40,opacity=0.55,line width=0.55pt]
                    ([xshift=-0.6cm,yshift=-0.4cm]current page.south west) circle[radius=\r cm];
            }
            % ---------- 右下螺旋线 ----------
            \begin{scope}[shift={([xshift=-2.9cm,yshift=3.2cm]current page.south east)}]
                \draw[booklight!45!bookteal,opacity=0.5,line width=0.55pt]
                    plot[domain=0:1260,samples=200,smooth,variable=\t]
                    ({0.0036*\t*cos(\t)},{0.0036*\t*sin(\t)});
            \end{scope}
            % ---------- 左上斜线组 ----------
            \foreach \i in {0,1,2}{
                \draw[bookteal!60!booklight,opacity=0.5,line width=0.6pt]
                    ([xshift={1.6cm+\i*0.45cm}]current page.north west) --
                    ([yshift={-1.6cm-\i*0.45cm}]current page.north west);
            }

            % ---------- 中下部视觉主体:外摆线(玫瑰形) ----------
            \begin{scope}[shift={([yshift=-16.6cm]current page.north)}]
                % 外圈
                \draw[booklight!35,opacity=0.6,line width=0.6pt]
                    (0,0) circle[radius=4.9cm];
                \draw[booklight!22,opacity=0.5,line width=0.5pt]
                    (0,0) circle[radius=5.25cm];
                % 主玫瑰线
                \draw[booklight!55!bookteal,opacity=0.8,line width=0.85pt]
                    plot[domain=0:1080,samples=400,smooth,variable=\t]
                    ({0.335*(8*cos(\t)-5*cos(2.6667*\t))},
                     {0.335*(8*sin(\t)-5*sin(2.6667*\t))});
                % 内层玫瑰线
                \draw[booklight!38!bookteal,opacity=0.6,line width=0.6pt]
                    plot[domain=0:1080,samples=400,smooth,variable=\t]
                    ({0.21*(8*cos(\t)-5*cos(2.6667*\t))},
                     {0.21*(8*sin(\t)-5*sin(2.6667*\t))});
            \end{scope}

            % ---------- 底部半透明波浪 ----------
            \fill[white,opacity=0.06]
                ([yshift=4.6cm]current page.south west)
                .. controls ([xshift=6.5cm,yshift=6.4cm]current page.south west)
                         and ([xshift=-6.0cm,yshift=2.6cm]current page.south east) ..
                ([yshift=5.2cm]current page.south east) --
                (current page.south east) -- (current page.south west) -- cycle;
            \fill[booklight,opacity=0.08]
                ([yshift=2.6cm]current page.south west)
                .. controls ([xshift=5.5cm,yshift=4.3cm]current page.south west)
                         and ([xshift=-6.5cm,yshift=1.0cm]current page.south east) ..
                ([yshift=3.3cm]current page.south east) --
                (current page.south east) -- (current page.south west) -- cycle;
        \end{tikzpicture}
        \mbox{}
    \end{titlepage}
    \endgroup
}
\makeatother

% 页脚右下方常态化水印(署名 + 仓库地址,加粗黑色),已在 main/plain 分页样式中定义,所有页面均生效.

\begin{document}
\mainmatter
\pagestyle{main}
\setcounter{chapter}{7}
\def\BookEnd{\end{document}}
\fi

% ==================== 附录:数学分析重要定理 ====================
% 本附录按小节组织,小节编号 §1,§2,§3,... 自动生成,
% 你只需写 \section{小节名},例如 \section{数列}.
% 子标题自动居左,编号前带段落符号 \S.
%
% 定理请放在 ti 环境(绿色框)中,证明写在框外下方,与正文题目结构一致:
%   \begin{ti}
%   【定理1】 定理陈述...
%   \end{ti}
%   \textbf{证明}:证明内容...
%
% 同样支持 \textbf{另解}、\textbf{注}、\textbf{注记}、\textbf{一点思考} 等标记.
% ================================================================

% 子标题编号格式:§1,§2,§3,...(带段落符号);标题居左.
\renewcommand{\thesection}{\S\arabic{section}}
\setcounter{section}{0}
\titleformat{\section}[block]
  {\normalfont\heiti\bfseries\Large\color{bookblue}}
  {\thesection}{1em}{}
\titlespacing*{\section}{0pt}{3.2em plus 0.4em}{1em}

% 子节(subsection)编号格式:§X.Y(父节号.子节号),标题居左,比小节略小.
\renewcommand{\thesubsection}{\S\arabic{section}.\arabic{subsection}}
\titleformat{\subsection}[block]
  {\normalfont\heiti\bfseries\large\color{bookblue}}
  {\thesubsection}{1em}{}
\titlespacing*{\subsection}{0pt}{2.4em plus 0.3em}{0.8em}

% 附录行间编号公式居中:正文已按 fleqn(左对齐)排版,本附录要求 equation/align 等编号公式居中,
% 故在此关闭 amsmath 的 fleqn 布局分支(无编号 $$...$$ 公式本就居中,不受影响)。仅作用于本附录及其后内容。
\makeatletter\@fleqnfalse\makeatother

% 附录公式编号:去掉"章.节"前缀,同一节内从 1 开始递推(1,2,3,...)。
\makeatletter
\@addtoreset{equation}{section}
\renewcommand{\theequation}{\arabic{equation}}
\makeatother

% 附录内局部放大绿框(ti)之间的间距:+20%(14pt → 16.8pt),不影响正文各章.
\renewtcolorbox{ti}{
    enhanced,
    breakable,
    colback=booklight!55!white,
    colframe=bookteal!35!white,
    boxrule=0.45pt,
    boxsep=0pt,
    left=10pt,
    right=10pt,
    top=9pt,
    bottom=9pt,
    before skip=16.8pt,
    after skip=16.8pt,
    arc=2pt,
    outer arc=2pt,
    before upper={\parindent=0pt}
}

\chapter*{附录\quad 数学分析重要定理}
\phantomsection
\addcontentsline{toc}{chapter}{附录\texorpdfstring{\quad}{ }数学分析重要定理}
\markboth{附录\quad 数学分析重要定理}{}

% 子标题与内容示例(由作者填写),例如:
% \section{数列}
% \begin{ti}
% 【定理1】 单调有界数列必收敛.
% \end{ti}
% \textbf{证明}: ...
%
% \section{级数}
% \begin{ti}
% 【定理1】 正项级数收敛的$Cauchy$判别法:...
% \end{ti}
% \textbf{证明}: ...

\section{数学分析前置知识}
\subsection{不等式}

\begin{ti}
\textbf{命题1.1(三角不等式)}  对任意的$a,b\in\mathbb{R}$,有$|a+b|\leq |a|+|b|$;
对任意的$a,b,c\in\mathbb{R}$,有$|a-b|\leq |a-c|+|c-b|$.
\end{ti}

\textbf{注}:三角不等式不仅对实数的绝对值成立,复数的模以及任意赋范向量空间中的范数也满足相应的三角不等式.

\begin{ti}
\textbf{定理1.2(算术平均值-几何平均值不等式)}  设$a_1,a_2,...,a_n$是$n$个非负实数,则成立不等式
$$\dfrac{a_1+a_2+...+a_n}{n}\geq\sqrt[n]{a_1a_2...a_n},$$
其中等号成立的充分必要条件为$a_1=a_2=...=a_n$.
\end{ti}

\textbf{证明}:关于n元平均值不等式的证明方法有很多,这里介绍用向前向后归纳法的证明
(这是由$Cauchy$给出的,同时在正文的第一章题3也用到了这种方法).\\

从$n=2$的已知情形出发,可以得到如下$n=4$时的平均值不等式:
$$\dfrac{a_1+a_2+a_3+a_4}{4}=\dfrac{1}{2}\left(\dfrac{a_1+a_2}{2}+\dfrac{a_3+a_4}{2}\right)$$
$$\phantom{\dfrac{a_1+a_2+a_3+a_4}{4}=}\geq\sqrt{\left(\dfrac{a_1+a_2}{2}\right)\cdot\left(\dfrac{a_3+a_4}{2}\right)}$$
$$\phantom{\dfrac{a_1+a_2+a_3+a_4}{4}=}\geq\sqrt{\sqrt{a_1a_2}\sqrt{a_3a_4}}=\sqrt[4]{a_1a_2a_3a_4}.$$
同样可知,若$n=2^k$时不等式成立,则可得到$n=2^{k+1}$时的平均值不等式
$$\dfrac{1}{2^{k+1}}\sum_{i=1}^{2^{k+1}}a_i=\dfrac{1}{2}\left(\dfrac{1}{2^k}\sum_{i=1}^{2^k}a_i+\dfrac{1}{2^k}\sum_{i=2^k+1}^{2^{k+1}}a_i\right)$$
$$\phantom{\dfrac{1}{2^{k+1}}\sum_{i=1}^{2^{k+1}}a_i=}\geq\sqrt{\left(\dfrac{1}{2^k}\sum_{i=1}^{2^k}a_i\right)\cdot\left(\dfrac{1}{2^k}\sum_{i=2^k+1}^{2^{k+1}}a_i\right)}$$
$$\phantom{\dfrac{1}{2^{k+1}}\sum_{i=1}^{2^{k+1}}a_i=\qquad}\geq\sqrt{\sqrt[2^k]{\prod_{i=1}^{2^k}a_i}\sqrt{\sqrt[2^k]{\prod_{i=2^k+1}^{2^{k+1}}a_i} } }=\sqrt[2^{k+1}]{\prod_{i=1}^{2^{k+1}}a_i}.$$
这样就证明了当$n$为2的所有方幂时平均值不等式已成立.这是向前部分($n\to 2n$).\\

第二步要证明,当平均值不等式对某个$n>2$成立时,则它对$n-1$也一定成立.
这是证明中的向后部分.写出
$$\dfrac{1}{n-1}\sum_{i=1}^{n-1}a_i=\dfrac{1}{n}\cdot\left(\dfrac{n}{n-1}\right)\sum_{i=1}^{n-1}a_i=\dfrac{1}{n}\left(\sum_{i=1}^{n-1}a_i+\dfrac{1}{n-1}\sum_{i=1}^{n-1}a_i\right).$$
将圆括号中的第二项看成为$a_n$,就可以利用$n$时已成立的平均值不等式得到
$$\dfrac{1}{n-1}\sum_{i=1}^{n-1}a_i\geq\sqrt[n]{\left(\prod_{i=1}^{n-1}a_i\right)\cdot\left(\dfrac{1}{n-1}\sum_{i=1}^{n-1}a_i\right)}.$$
将以上不等式两边升高$n$次幂,就有
$$\left(\dfrac{1}{n-1}\sum_{i=1}^{n-1}a_i\right)^n\geq\left(\prod_{i=1}^{n-1}a_i\right)\cdot\left(\dfrac{1}{n-1}\sum_{i=1}^{n-1}a_i\right),$$
若$a_1=\cdots=a_{n-1}=0$,所要的不等式显然成立;否则上述平均值为正,可以在两边约去公因子$\displaystyle\dfrac{1}{n-1}\sum_{i=1}^{n-1}a_i$,再开$n-1$次根,就得到所要的不等式.\\
合并以上向前向后的两部分,可见平均值不等式对每个正整数$n$都成立.

\begin{ti}
\textbf{命题1.3($Newton$二项式展开)}  设$n$为正整数,则对任意的$a,b\in\mathbb{R}$,有
$$(a+b)^{n}=\sum_{k=0}^{n}\binom{n}{k}a^{k}b^{n-k}.$$
\end{ti}

\begin{ti}
\textbf{定理1.4($Bernoulli$不等式)} 设$h>-1,n\in\mathbb{N}^+,$则成立不等式
$$(1+h)^n\geq 1+nh,$$
其中当$n>1$时等号成立的充分必要条件为$h=0$.
\end{ti}

\begin{ti}
\textbf{命题1.5($Cauchy$不等式)} 对实数$a_1,a_2,...,a_n$和$b_1,b_2,...,b_n$成立
$$\left\lvert \sum_{i=1}^{n}a_ib_i\right\rvert\leq\sqrt{\sum_{i=1}^{n}a_i^2}\sqrt{\sum_{i=1}^{n}b_i^2}.$$
\end{ti}

\subsection{欧几里得空间}
\begin{ti}
\textbf{定义1.6(欧氏距离、开球与邻域)}  对$x=(x_1,\ldots,x_n),y=(y_1,\ldots,y_n)\in\mathbb{R}^n$,定义欧氏范数$\|x\|=\sqrt{\sum_{i=1}^{n}x_i^2}$和欧氏距离$d(x,y)=\|x-y\|$.给定$r>0$,以$x$为中心、$r$为半径的\textbf{开球}为
$$B_r(x)=\{y\in\mathbb{R}^n:d(x,y)<r\}.$$
若集合$V\subset\mathbb{R}^n$包含某个$B_r(x)$,则称$V$为点$x$的一个\textbf{邻域}.当$n=1$时,$B_r(x)=(x-r,x+r)$.
\end{ti}

\begin{ti}
\textbf{定理1.7(确界原理)}  设$A\subset\mathbb{R}$为非空子集.如果$A$有上界,则它有一个最小上界,称为$A$的\textbf{上确界},记为$\sup A$;
如果$A$有下界,则它有一个最大下界,称为$A$的\textbf{下确界},记为$\inf A$.
\end{ti}

\begin{ti}
\textbf{命题1.8(确界的刻画)}  设$A\subset\mathbb{R}$为非空子集.
如果$A$以$\beta$为上界,则$\beta$为$A$的上确界当且仅当任给正整数$n$,均存在$a_n\in A$,使得$a_n>\beta-\dfrac{1}{n}$;
如果$A$以$\alpha$为下界,则$\alpha$为$A$的下确界当且仅当任给正整数$n$,均存在$a_n\in A$,使得$a_n<\alpha+\dfrac{1}{n}$.\\

\textbf{注}:与之等价的一种刻画为:$\beta$为$A$的上确界当且仅当任给$\varepsilon>0$,均存在$a\in A$,使得$a>\beta-\varepsilon$;
$\alpha$为$A$的下确界当且仅当任给$\varepsilon>0$,均存在$b\in A$,使得$b<\alpha+\varepsilon$.
\end{ti}

\begin{ti}
\textbf{定理1.9(阿基米德原理)}  设$x\in\mathbb{R}^{+}$,则任给$y\in\mathbb{R}$,存在正整数$n$,使得$y<nx$.
\end{ti}

\begin{ti}
\textbf{定理1.10(有理数的稠密性)}  任给$a,b\in\mathbb{R}$,如果$a<b$,则存在$c\in\mathbb{Q}$,使得$a<c<b$.
\end{ti}

\subsection{集合与映射}
\begin{ti}
\textbf{定义1.11(集合运算)}  设$A,B$为固定全集$X$的子集.并集$A\cup B=\{x\in X:x\in A\text{或}x\in B\}$,交集$A\cap B=\{x\in X:x\in A\text{且}x\in B\}$,补集$A^c=X\setminus A$.对任意指标集$I$和由$X$的子集组成的集合族$\{A_i\}_{i\in I}$,有$De\,Morgan$律
$$\left(\bigcup_{i\in I}A_i\right)^c=\bigcap_{i\in I}A_i^c,\qquad
\left(\bigcap_{i\in I}A_i\right)^c=\bigcup_{i\in I}A_i^c.$$
补集均相对于同一个全集$X$而言.
\end{ti}

\begin{ti}
\textbf{定义1.12(开集与闭集)}  设$U,F\subset\mathbb{R}^n$.若对每个$x\in U$,都存在$r>0$使$B_r(x)\subset U$,则称$U$为\textbf{开集};若$\mathbb{R}^n\setminus F$为开集,则称$F$为\textbf{闭集}.空集与$\mathbb{R}^n$既是开集也是闭集.
\end{ti}

\begin{ti}
\textbf{命题1.13(开闭集的并与交)}  在$\mathbb{R}^n$中,任意个开集的并、有限个开集的交仍为开集;任意个闭集的交、有限个闭集的并仍为闭集.
\end{ti}

\textbf{证明}:若$x$属于一族开集的并,则$x$属于其中某个开集,该开集包含以$x$为中心的一个开球,故并集为开集.若$x$属于有限个开集的交,分别取各开集中以$x$为中心的开球,以这些开球半径的最小值为半径的开球包含于交集中,故交集为开集.闭集的两条结论由补集的定义及$De\,Morgan$律得到.

\textbf{注}:上述“有限”不能任意去掉.在$\mathbb{R}$中,$\displaystyle\bigcap_{k=1}^{\infty}(-1/k,1/k)=\{0\}$不是开集,而$\displaystyle\bigcup_{k=2}^{\infty}[0,1-1/k]=[0,1)$不是闭集.

\begin{ti}
\textbf{定义1.14(相对开集与相对闭集)}  设$Y\subset\mathbb{R}^n$,$A\subset Y$.若存在$\mathbb{R}^n$中的开集$U$使$A=Y\cap U$,则称$A$在$Y$中\textbf{相对开};若存在$\mathbb{R}^n$中的闭集$F$使$A=Y\cap F$,则称$A$在$Y$中\textbf{相对闭}.例如$[0,1)$在$[0,2]$中相对开,但在$\mathbb{R}$中不是开集.
\end{ti}

\begin{ti}
\textbf{定义1.15(紧集)}  设$K\subset\mathbb{R}^n$.若对每一族满足$K\subset\bigcup_{i\in I}U_i$的开集$\{U_i\}_{i\in I}$,均存在有限子集$J\subset I$,使$K\subset\bigcup_{i\in J}U_i$,则称$K$为\textbf{紧集}.这里的$\{U_i\}$称为$K$的开覆盖.
\end{ti}

\begin{ti}
\textbf{定理1.16($Heine$--$Borel$定理)}  $K\subset\mathbb{R}^n$是紧集当且仅当$K$既闭又有界.这里的“有界”是指存在$R>0$,使$K\subset B_R(0)$.此等价关系依赖于欧几里得空间,不能直接推广到任意度量空间.
\end{ti}

\textbf{证明}:先设$K$紧.开球族$\{B_m(0)\}_{m\geq1}$覆盖$K$,由有限子覆盖知$K$有界.若$K\neq\varnothing$且$z\notin K$,对每个$x\in K$取$r_x=\|x-z\|/3$,则$\{B_{r_x}(x)\}_{x\in K}$覆盖$K$.取有限子覆盖,并令$r$为其中各$r_x$的最小值,便有$B_r(z)\cap K=\varnothing$,故$K$闭;空集显然闭.\\

反之,设$K$闭且有界,可将$K$置于一个闭立方体$Q_0$中.若某开覆盖不能有限覆盖$K$,则将$Q_0$沿各坐标方向二等分,必有一个子立方体与$K$的交仍不能被有限覆盖.如此逐次选取,得到套叠闭立方体$Q_j$,其直径趋于$0$,且每个$K\cap Q_j$都不能被有限覆盖.这些集合非空;由坐标分别应用闭区间套定理,有唯一$x\in\bigcap_jQ_j$.因为$K$闭,且可从每个$K\cap Q_j$选点趋于$x$,所以$x\in K$.开覆盖中某个开集包含$x$,因$Q_j$的直径趋于$0$,充分大的$Q_j$包含于这个开集,与$K\cap Q_j$不能被有限覆盖矛盾.

\begin{ti}
\textbf{命题1.17(紧集的基本性质)}  $\mathbb{R}^n$中紧集的闭子集仍紧;有限个紧集的并仍紧.
\end{ti}

\textbf{证明}:若$F\subset K$在$\mathbb{R}^n$中闭且$K$紧,则对$F$的任意开覆盖,再添上开集$\mathbb{R}^n\setminus F$,便得到$K$的开覆盖.取其有限子覆盖后去掉$\mathbb{R}^n\setminus F$,余下各开集仍覆盖$F$.对于有限个紧集的并,分别从同一个开覆盖中为各紧集选取有限子覆盖,合在一起即得所需的有限子覆盖.

\begin{ti}
\textbf{定义1.18(映射,单射,满射与双射)}  设$X,Y$为集合.
若对每个$x\in X$,均存在唯一的$y\in Y$与$x$相对应,则称$f$为从$X$到$Y$的\textbf{映射},记作$f:X\to Y$,并记为$x\mapsto f(x)$.设$f:X\to Y$为映射.\\
(1)如果任给$x_1,x_2\in X$,$f(x_1)=f(x_2)$时必有$x_1=x_2$,则称$f$为\textbf{单射};\\
(2)如果任给$y\in Y$,均存在$x\in X$,使得$f(x)=y$,则称$f$为\textbf{满射};\\
(3)如果$f$既是单射又是满射,则称$f$为\textbf{一一映射}或\textbf{双射}.
\end{ti}

\begin{ti}
\textbf{定义1.19(可数集)}  如果存在从正整数集$\mathbb{N}^+$到集合$A$之间的一一映射,则称$A$为\textbf{可数集}.
直观地看,集合$A$为可数集当且仅当其元素可以不重复也不遗漏地排成一列,即$A$可以表示为$A=\{a_1,a_2,\cdots,a_n,\cdots\}$.只含有有限个元素的集合称为\textbf{有限集};
有限集和可数集统称\textbf{至多可数集},其余集合称为\textbf{不可数集}.
\end{ti}

\begin{ti}
\textbf{命题1.20}  有理数集$\mathbb{Q}$为可数集.
\end{ti}

\textbf{证明}:对每一对正整数$(p,q)$,按$p+q$从小到大排列;在$p+q$相同时,按$p$从小到大排列.依此顺序写出$\displaystyle\frac{p}{q}$,但遇到$(p,q)>1$的分数就略去.所得正有理数列以
$$1,\frac{1}{2},2,\frac{1}{3},3,\frac{1}{4},\frac{2}{3},\frac{3}{2},4,\cdots$$
开头,记为$r_1,r_2,\ldots$.每个正有理数都能唯一地写成$\displaystyle\frac{p}{q}$,其中$p,q\in\mathbb{N}^+$且$\gcd(p,q)=1$;该数在处理到$p+q$时必被写出,而互素分子分母的表示唯一,所以此列无遗漏、无重复.最后按
$$0,r_1,-r_1,r_2,-r_2,\ldots$$
排列,零仅出现一次,而每个非零有理数的绝对值恰为某个唯一的$r_i$,其符号也确定了它在此列中的位置.因此这一排列无重复、无遗漏,即为全体有理数的一一枚举,故$\mathbb{Q}$可数.

\begin{ti}
\textbf{命题1.21}  实数集$\mathbb{R}$为不可数集.
\end{ti}

\textbf{证明}:用$Cantor$对角线法.反设$(0,1)$中的实数可以排成一列$x_1,x_2,\ldots$.将每个$x_n$写成不以无限重复的$9$结尾的十进制小数(如此表示唯一)
$$x_n=0.d_{n1}d_{n2}d_{n3}\cdots\qquad(n=1,2,\ldots).$$
逐位构造$y=0.b_1b_2b_3\cdots$:若$d_{nn}=1$,则取$b_n=2$;否则取$b_n=1$.由于$y$的小数位全为$1$或$2$,它属于$(0,1)$,且其表示也不以无限重复的$9$结尾.对每个$n$,数$y$与$x_n$的第$n$位不同,故$y\ne x_n$.这与$x_1,x_2,\ldots$已列出$(0,1)$中的所有实数矛盾.因此$(0,1)$不可数;若$\mathbb{R}$可数,从其枚举中删去不属于$(0,1)$的数便能枚举$(0,1)$,仍得矛盾.

\section{数列}
\begin{ti}
\textbf{定义2.1(数列极限)}  设$\{a_n\}$为数列,$\alpha\in\mathbb{R}$,如果任给$\varepsilon>0$,均存在正整数$N=N(\varepsilon)$,
使得当$n>N$时,$\left\lvert a_n-\alpha\right\rvert<\varepsilon$总成立,则称$\{a_n\}$收敛于$\alpha$,或称$\{a_n\}$以$\alpha$为极限,记为
$$\displaystyle\lim_{n\to\infty}a_n=\alpha\quad\text{或}\quad a_n\rightarrow\alpha$$.
\end{ti}

\begin{ti}
\textbf{定理2.2(单调数列的极限)}  设$\{a_n\}$为单调数列. \\

(1)如果$\{a_n\}$单调递增且有上界,则$\displaystyle\lim_{n\to\infty}a_n=\sup\left\{a_k\;|\;k\geq 1\right\}$;\\
(2)如果$\{a_n\}$单调递减且有下界,则$\displaystyle\lim_{n\to\infty}a_n=\inf\left\{a_k\;|\;k\geq 1\right\}$.
\end{ti}

\begin{ti}
\textbf{推论2.3} 有界单调数列必收敛
\end{ti}

\begin{ti}
\textbf{定义2.4(上下极限)} 对于有界数列$\{a_n\}$,令
$$\underline{a_n}=\inf\left\{a_k\;|\;k\geq n\right\},\quad \overline{a_n}=\sup\left\{a_k\;|\;k\geq n\right\}$$.
易见$\{\underline{a_n}\}$和$\{\overline{a_n}\}$分别是单调递增和单调递减的数列,且
$$\underline{a_n}\leq a_n\leq\overline{a_n}$$.
单调数列$\{\underline{a_n}\}$和$\{\overline{a_n}\}$的极限称为$\{a_n\}$的\textbf{下极限}和\textbf{上极限},记为
$$\displaystyle\varliminf_{n\to\infty}a_n=\lim_{n\to\infty}\underline{a_n},\quad\varlimsup_{n\to\infty}a_n=\lim_{n\to\infty}\overline{a_n}$$.
\end{ti}

\begin{ti}
\textbf{定理2.5} 设$\{a_n\}$为有界数列,则下列命题等价: \\

(1)$\{a_n\}$收敛;\\
(2)$\{a_n\}$的上极限和下极限相等;\\
(3)$\displaystyle\lim_{n\to\infty}(\overline{a_n}-\underline{a_n})=0$.
\end{ti}

\begin{ti}
\textbf{命题2.6} 设$\{a_n\},\{b_n\}$为有界数列.\\
(1)(上下极限的保序性质)如果存在$N_0,n>N_0$时$a_n\geq b_n$,则
$$\displaystyle\varliminf_{n\to\infty}a_n\geq\varliminf_{n\to\infty}b_n,\quad\varlimsup_{n\to\infty}a_n\geq\varlimsup_{n\to\infty}b_n$$;
(2)$\displaystyle\varlimsup_{n\to\infty}(a_n+b_n)\leq\varlimsup_{n\to\infty}a_n+\varlimsup_{n\to\infty}b_n$.
\end{ti}

\begin{ti}
\textbf{定理2.7($Cauchy$准则)} 数列$\{a_n\}$收敛当且仅当它是$Cauchy$数列.
\end{ti}

\textbf{证明}:必要性:设$\{a_n\}$收敛到$\alpha$.任给$\varepsilon>0$,存在$N$,当$n>N$时$\left\lvert a_n-\alpha\right\rvert<\dfrac{\varepsilon}{2}$.因此,当$m,n>N$时,有
$$\left\lvert a_m-a_n\right\rvert\leq\left\lvert a_m-\alpha\right\rvert+\left\lvert\alpha-a_n\right\rvert<\frac{\varepsilon}{2}+\frac{\varepsilon}{2}=\varepsilon,$$
这说明$\{a_n\}$为$Cauchy$数列. \\

\phantom{\textbf{证明}:}充分性:设$\{a_n\}$为$Cauchy$数列.由于$Cauchy$数列必有界,$\{a_n\}$是有界数列,于是可以研究其上、下极限.根据$Cauchy$数列的定义,
任给$\varepsilon>0$,存在$N$,当$m,n>N$时,
$$-\varepsilon<a_m-a_n<\varepsilon.$$
在上式中暂时固定$n>N$,对$\{a_m\}$取上极限,利用上极限的保序性可得
$$-\varepsilon\leq\displaystyle\varlimsup_{m\to\infty}a_m-a_n\leq\varepsilon.$$
由数列极限的定义即可看出$\{a_n\}$收敛.

\begin{ti}
\textbf{定理2.8($Stolz$公式)} 设$\{x_n\},\{y_n\}$为数列, \\

(1)若$\{y_n\}$严格单调地趋于$+\infty,\displaystyle\lim_{n\to\infty}\dfrac{x_n-x_{n-1}}{y_n-y_{n-1}}=\alpha,\alpha$为实数或$\pm\infty$,则
$$\displaystyle\lim_{n\to\infty}\dfrac{x_n}{y_n}=\alpha$$. \\
(2)若$\{y_n\}$严格单调递减趋于$0,\{x_n\}$也收敛到$0,\displaystyle\lim_{n\to\infty}\dfrac{x_n-x_{n-1}}{y_n-y_{n-1}}=\alpha,\alpha$为实数或$\pm\infty$,则
$$\displaystyle\lim_{n\to\infty}\dfrac{x_n}{y_n}=\alpha$$.
\end{ti}

\textbf{证明}:先证(1).分情况讨论. \\
\romannum{1}$\;\alpha\in\mathbb{R}$.先设$\alpha=0$.此时,任给$\varepsilon>0$,存在$N$,当$n>N$时,
$$-\varepsilon<\dfrac{x_n-x_{n-1}}{y_n-y_{n-1}}<\varepsilon.$$
利用$\{y_n\}$的单调性以及引理(当$1\leq k\leq n$时,设$b_k>0$且$m\leq\dfrac{a_k}{b_k}\leq M$,则$m\leq\dfrac{a_1+\cdots+a_n}{b_1+\cdots+b_n}\leq M$),当$n>N$时,得到
$$\left\lvert\dfrac{x_n-x_N}{y_n-y_N}\right\rvert<\varepsilon,$$
整理以后可得
$$\frac{x_N+\varepsilon y_N}{y_n}-\varepsilon<\frac{x_n}{y_n}<\frac{x_N-\varepsilon y_N}{y_n}+\varepsilon.$$
由$\displaystyle\lim_{n\to\infty}y_n=+\infty$以及上式可知,当$n$充分大时,$-2\varepsilon<\dfrac{x_n}{y_n}<2\varepsilon$,因此$\displaystyle\lim_{n\to\infty}\dfrac{x_n}{y_n}=0$. \\
一般地,设$\displaystyle\lim_{n\to\infty}\dfrac{x_n-x_{n-1}}{y_n-y_{n-1}}=\alpha$,记$\widetilde{x_n}=x_n-\alpha y_n$,则
$$\lim_{n\to\infty}\dfrac{\widetilde{x_n}-\widetilde{x_{n-1}}}{y_n-y_{n-1}}=\lim_{n\to\infty}\left(\dfrac{x_n-x_{n-1}}{y_n-y_{n-1}}-\alpha\right)=0,$$
这说明$\displaystyle\lim_{n\to\infty}\dfrac{\widetilde{x_n}}{y_n}=0$,从而$\displaystyle\lim_{n\to\infty}\dfrac{x_n}{y_n}=\alpha$. \\

\romannum{2}$\;\alpha=+\infty$.此时,存在$N$,当$n>N$时,$\dfrac{x_n-x_{n-1}}{y_n-y_{n-1}}>1$.于是,当$n>N$时,
$$x_n-x_{n-1}>y_n-y_{n-1}>0,$$
即$n>N$时$\{x_n\}$也是严格单调递增的,且
$$x_n-x_N=(x_n-x_{n-1})+\cdots+(x_{N+1}-x_N)>(y_n-y_N),$$
由$\displaystyle\lim_{n\to\infty}y_n=+\infty$可知$\displaystyle\lim_{n\to\infty}x_n=+\infty$.
由题设可得$\displaystyle\lim_{n\to\infty}\dfrac{y_n-y_{n-1}}{x_n-x_{n-1}}=0$.
应用(1)的结论可得$\displaystyle\lim_{n\to\infty}\dfrac{y_n}{x_n}=0$,于是$\displaystyle\lim_{n\to\infty}\dfrac{x_n}{y_n}=+\infty$. \\

\romannum{3}$\;\alpha=-\infty$.这时只要将$x_n$换成$-x_n$,然后应用本证明中$\alpha=+\infty$的情形即可. \\

再证(2). \\
\romannum{1}$\;\alpha\in\mathbb{R}$.先设$\alpha=0$.任给$\varepsilon>0$,存在$N$,当$n>N$时,$\left\lvert\dfrac{x_{n+1}-x_n}{y_{n+1}-y_n}\right\rvert<\varepsilon$,和上面的证明类似,当$m>n>N$时可得
$$-\varepsilon(y_n-y_m)<x_n-x_m<\varepsilon(y_n-y_m),$$
令$m\to\infty$,可得$-\varepsilon y_n\leq x_n\leq\varepsilon y_n$,这说明$\displaystyle\lim_{n\to\infty}\dfrac{x_n}{y_n}=0$.一般的有限$\alpha$只需将$x_n$替换为$x_n-\alpha y_n$. \\

\romannum{2}$\;\alpha=+\infty$.任给$M>0$,存在$N$,当$n>N$时,$\dfrac{x_n-x_{n+1}}{y_n-y_{n+1}}>M$.
与上面的证明类似,当$m>n>N$时,有$x_n-x_m>M(y_n-y_m)$,令$m\to\infty$,得$x_n\geq My_n$,这说明$\displaystyle\lim_{n\to\infty}\dfrac{x_n}{y_n}=+\infty$. \\

\romannum{3}$\;\alpha=-\infty$.将(2)中$x_n$换成$-x_n$即可.

\begin{ti}
\textbf{定理2.9($Wallis$公式)}  $\displaystyle\lim_{n\to\infty}\dfrac{1}{2n+1}\left[\dfrac{(2n)!!}{(2n-1)!!}\right]^{2}=\dfrac{\pi}{2}.$
\end{ti}
\textbf{证明}:瓦利斯公式的证明需要用到三角函数的积分知识,详细证明请见定理5.19
\begin{ti}
\textbf{定理2.10($Stirling$公式)} $n!=\sqrt{2\pi n}\left(\dfrac{n}{e}\right)^n e^{\delta_n},\;0<\delta_n<\dfrac{1}{12n}$.
\end{ti}
\textbf{证明}:斯特林公式的证明需要用到积分学和泰勒展开的知识,详细证明请见定理5.20.

\section{实数系基本定理}
\begin{ti}
\textbf{定理3.1($Cantor$)} 设$\{[a_n,b_n]\}$是闭区间套,如果$\displaystyle\lim_{n\to\infty}(b_n-a_n)=0,$则这些闭区间有一个唯一的公共点.
\end{ti}

\textbf{证明}:由题设可知,$\{a_n\}$单调递增,$\{b_n\}$单调递减,它们均位于区间$[a_1,b_1]$中,从而为有界数列.
由单调有界数列必收敛,$\{a_n\}$和$\{b_n\}$都收敛,设其极限分别为$\alpha,\beta$.
由$a_n<b_n$和极限的保序性质可知$\alpha\leq\beta$.由$\{a_n\}$和$\{b_n\}$的单调性可知$a_n\leq\alpha,\beta\leq b_n$.
于是$\alpha,\beta$为$\{[a_n,b_n]\}$的公共点,且
$$0\leq\beta-\alpha\leq b_n-a_n,\;\forall\,n\geq 1.$$
令$n\to\infty$,由$\displaystyle\lim_{n\to\infty}(b_n-a_n)=0$即得$\alpha=\beta$.
如果$\{[a_n,b_n]\}$另有公共点$\gamma$,则
$$0\leq\left\lvert\gamma-\alpha\right\rvert\leq b_n-a_n,\;\forall\,n\geq 1,$$同理可得$\gamma=\alpha$.

\begin{ti}
\textbf{定理3.2($Heine$--$Borel$)} 设$\{U_{\alpha}\}_{\alpha\in\varGamma}$为$\mathbb{R}$中的一族开集,如果闭区间$[a,b]$包含于这一族开集的并集之中,则$[a,b]$
必定包含于有限个$U_{\alpha}$的并集之中.
\end{ti}

\textbf{证明}:如果指标集$\varGamma$是有限集,则结论不证自明.以下设$\varGamma$为无限集.\\
(反证法)假设$[a,b]$只能包含于无限个$U_{\alpha}$的并集,则二等分$[a,b]$后必有一个小区间也只能包含于无限个$U_{\alpha}$之并,记该区间为$[a_1,b_1]$.
再将$[a_1,b_1]$二等分,又必有一个小区间只能包含于无限个$U_{\alpha}$之并,记该区间为$[a_2,b_2]$.
如此继续下去,得闭区间套$\{[a_n,b_n]\}$,使得每一个$[a_n,b_n]$均只能包含于无限个$U_{\alpha}$之并.\\

注意$\displaystyle\lim_{n\to\infty}(b_n-a_n)=\lim_{n\to\infty}2^{-n}(b-a)=0$.根据闭区间套原理,$\{[a_n,b_n]\}$有一个公共点,记为$\xi$.
根据题设,存在$\alpha_0\in\varGamma$,使得$\xi\in U_{\alpha_0}$.因为$U_{\alpha_0}$是开集,故存在$\delta>0$,使得$(\xi-\delta,\xi+\delta)\subset U_{\alpha_0}$.
根据闭区间套原理的证明,$\{a_n\},\{b_n\}$均收敛于$\xi$,故存在$N$,当$n>N$时$a_n,b_n\in(\xi-\delta,\xi+\delta)$.此时有
$$[a_n,b_n]\subset(\xi-\delta,\xi+\delta)\subset U_{\alpha_0},$$
这与$[a_n,b_n]$只能包含于无限个$U_{\alpha}$之并相矛盾.

\begin{ti}
\textbf{定理3.3($Bolzano$)}有界数列必有收敛子列
\end{ti}

\textbf{证明}:设$\{a_n\}$为有界数列,不妨设$a_n$均包含于$[a,b]$.\\
\phantom{\textbf{证明}:}断言:存在$\alpha\in[a,b]$,使得任给$\delta>0$,$(\alpha-\delta,\alpha+\delta)$中均含有无限项$a_n$.\\
\phantom{\textbf{证明}:}(反证法)假设不然,则任给$x\in[a,b]$,存在$\delta(x)>0$,使得$(x-\delta(x),x+\delta(x))$只含有限项$a_n$.
显然,$[a,b]$包含于集合族$\{(x-\delta(x),x+\delta(x))\}_{x\in[a,b]}$之并.由$Heine$--$Borel$定理,$[a,b]$包含于有限个$(x-\delta(x),x+\delta(x))$之并.
这说明$[a,b]$只含有限项$a_n$,与$a_n$均包含于$[a,b]$相矛盾.利用此断言,我们选取$\{a_n\}$的子列,使之收敛到$\alpha$.
事实上,先取$a_{n_1}\in(\alpha-1,\alpha+1)$.再取$n_2>n_1$,使得$a_{n_2}\in\left(\alpha-\dfrac{1}{2},\alpha+\dfrac{1}{2}\right)$.
如此继续,可得子列$\{a_{n_k}\}$,使得当$k\geq 1$时$a_{n_k}\in\left(\alpha-\dfrac{1}{k},\alpha+\dfrac{1}{k}\right)$.显然,$\{a_{n_k}\}$收敛到$\alpha$.

\section{函数}
\begin{ti}
\textbf{定义4.1(函数极限)} 设函数$f$在$x_0$的一个空心开邻域$V(x_0,\delta_0)$中有定义.如果存在$\alpha\in\mathbb{R}$,使得任给$\varepsilon>0$,
均存在$\delta\in(0,\delta_0)$,当$x\in V(x_0,\delta)$时,$\left\lvert f(x)-\alpha\right\rvert<\varepsilon$,则称函数$f$在$x_0$处有极限$\alpha$,记为
$$\displaystyle\lim_{x\to x_0}f(x)=\alpha\;\text{或}\;f(x)\rightarrow\alpha(x\rightarrow x_0)$$.
\end{ti}

\begin{ti}
\textbf{命题4.2} \\
(1)(局部有界性)如果$f$在$x_0$处有极限,则它在$x_0$的一个空心开邻域中有界.\\
(2)(保序性)设函数$f,g$在$x_0$处的极限分别为$\alpha,\beta$.如果$f(x)\geq g(x)$在$x_0$的一个空心开邻域中成立,则$\alpha\geq\beta$,
反之,如果$\alpha>\beta$,则在$x_0$的一个空心开邻域中$f(x)>g(x)$;特别地,如果$\alpha>0$,则在$x_0$的一个空心开邻域中$f(x)>0$
\end{ti}

\begin{ti}
\textbf{定理4.3(复合函数的极限)}设函数$f(y)$在$y_0$处极限为$\alpha$,函数$g(x)$在$x_0$处极限为$y_0$,
如果存在$x_0$的一个空心开邻域,在此开邻域中$g(x)\neq y_0$,则复合函数$f(g(x))$在$x_0$处极限为$\alpha$
\end{ti}

\begin{ti}
\textbf{定理4.4($Heine$)} 设函数$f$在$x_0$的一个空心开邻域中有定义,则$f$在$x_0$
处极限为$\alpha$当且仅当对空心开邻域中任何收敛于$x_0$的数列$\{x_n\}$,均有
$\displaystyle\lim_{n\to\infty}f(x_n)=\alpha$.
\end{ti}

\begin{ti}
\textbf{定义4.5(无穷小量与无穷大量)}  设$x_0$为实数或$\pm\infty$.\\
(1)如果$\displaystyle\lim_{x\to x_0}f(x)=0$,则称函数$f$在$x\to x_0$时为\textbf{无穷小量},记为$f(x)=o(1)\,(x\to x_0)$;\\
(2)如果$\dfrac{1}{f}$在$x\to x_0$时为无穷小量,则称$f$在$x\to x_0$时为\textbf{无穷大量}.
\end{ti}

\begin{ti}
\textbf{命题4.6(无穷小量的比较)}  设$f,g$均为$x\to x_0$时的无穷小量.\\
(1)如果$\dfrac{f(x)}{g(x)}=o(1)\,(x\to x_0)$,则称当$x\to x_0$时$f$关于$g$为\textbf{高(低)阶}无穷小量,记为$f(x)=o(g(x))\,(x\to x_0)$;\\
(2)如果$\dfrac{f(x)}{g(x)}$在$x_0$处的极限为非零实数,则称当$x\to x_0$时$f$与$g$为\textbf{同阶量};\\
(3)如果$\dfrac{f(x)}{g(x)}$在$x_0$处的极限等于$1$,则称当$x\to x_0$时$f$与$g$为\textbf{等价量},记为$f\sim g\,(x\to x_0)$.
\end{ti}

\begin{ti}
\textbf{定义4.7(函数的连续性)}  设函数$f$在$x_0$的某个邻域中有定义.如果$\displaystyle\lim_{x\to x_0}f(x)=f(x_0)$,则称$f$在$x_0$处\textbf{连续}.
如果$f$在区间$I$的每一点处均连续,则称$f$在$I$上连续.
\end{ti}

\begin{ti}
\textbf{定义4.8(间断点)}  设函数$f$在$x_0$的某个空心开邻域中有定义.如果$f$在$x_0$处不连续,则称$x_0$为$f$的\textbf{间断点}.
若$f$在$x_0$处的左右极限均存在,则称$x_0$为\textbf{第一类间断点},其中左右极限不等者称为\textbf{跳跃间断点},左右极限相等者称为\textbf{可去间断点};
其余情形称为\textbf{第二类间断点}.
\end{ti}

\begin{ti}
\textbf{定理4.9(最值定理)}  设函数$f$在闭区间$[a,b]$上连续,则$f$在$[a,b]$中必定取到最大值和最小值,即存在$x_1,x_2\in[a,b]$,使得
$$f(x_1)\leq f(x)\leq f(x_2),\quad \forall x\in[a,b].$$
\end{ti}

\textbf{证明}:由连续函数的有界性质,$f$有界,因此必有上确界和下确界.记上确界为$M$.
根据确界的刻画,存在$[a,b]$中的数列$\{x_n\}$,使得$M-\frac{1}{n}\leq f(x_n)\leq M$.根据$Bolzano$定理,$\{x_n\}$有收敛子列$\{x_{n_i}\}$,其极限记为$x_*$.
由极限的保序性质可知$x_*\in[a,b]$.因为$f$在$x_*$处连续,故$\{f(x_{n_i})\}$收敛于$f(x_*)$.这说明$f(x_*)=M$,$M$即为$f$的最大值.同理可证$f$可以取到最小值(或考虑$-f$的最大值).

\begin{ti}
\textbf{定理4.10(介值定理)}  设函数$f$在闭区间$[a,b]$上连续,$\mu$是严格介于$f(a)$与$f(b)$之间的数,则存在$\xi\in(a,b)$,使得$f(\xi)=\mu$.
\end{ti}

\textbf{证明}:不妨设$f(a)<\mu<f(b)$,并令$E=\{x\in[a,b]:f(x)<\mu\}$.因$a\in E$,集合$E$非空且有上界,故$\xi=\sup E$存在.由$f$在$a$处连续且$f(a)<\mu$,存在$a$右侧的点属于$E$,所以$\xi>a$;由$f$在$b$处连续且$f(b)>\mu$,存在$b$左侧的一段区间与$E$不交,所以$\xi<b$.若$f(\xi)<\mu$,连续性使$\xi$右侧某点仍属于$E$,与上确界的定义矛盾;若$f(\xi)>\mu$,连续性使$\xi$左侧某段区间与$E$不交,也与上确界的定义矛盾.故$f(\xi)=\mu$.当$f(a)>\mu>f(b)$时,对$-f$应用上述结论即可.

\begin{ti}
\textbf{推论4.11(零值定理,$Bolzano$)}  设函数$f$在闭区间$[a,b]$上连续,且$f(a)f(b)\leq 0$,则存在$\xi\in[a,b]$,使得$f(\xi)=0$.
\end{ti}

\textbf{证明}:若$f(a)f(b)=0$,取零值端点即可;否则$0$严格介于$f(a)$与$f(b)$之间,由定理4.10即得.

\begin{ti}
\textbf{定义4.12(一致连续)}  设函数$f$在区间$I$上有定义.如果对任意的$\varepsilon>0$,均存在$\delta>0$,使得对任意的$x_1,x_2\in I$,
只要$\left\lvert x_1-x_2\right\rvert<\delta$,就有$\left\lvert f(x_1)-f(x_2)\right\rvert<\varepsilon$,则称$f$在$I$上\textbf{一致连续}.
\end{ti}

\begin{ti}
\textbf{定理4.13($Cantor$)}  设函数$f$在闭区间$[a,b]$上连续,则$f$在$[a,b]$上一致连续.
\end{ti}

\textbf{证明}:设$f\in C^0[a,b]$.(反证法)若$f$不一致连续,则存在$\varepsilon_0>0$,以及$[a,b]$中的数列$\{a_n\},\{b_n\}$,使得
$$\lim_{n\to\infty}(a_n-b_n)=0,\qquad \left\lvert f(a_n)-f(b_n)\right\rvert\geq\varepsilon_0,\;\forall n\geq 1.$$
\phantom{\textbf{证明}:}根据$Bolzano$定理,$\{b_n\}$有收敛子列$\{b_{n_i}\}$,设其极限为$\xi$,
则$\xi\in[a,b]$.此时$a_{n_i}=(a_{n_i}-b_{n_i})+b_{n_i}\to 0+\xi=\xi\;(i\to\infty)$.因为$f$在$\xi$处连续,故
$$\varepsilon_0\leq\left\lvert f(a_{n_i})-f(b_{n_i})\right\rvert\to\left\lvert f(\xi)-f(\xi)\right\rvert=0\;(i\to\infty),$$
这就导出了矛盾.


\BookEnd
